\section{Limitations}\label{sec:limitations}
The service is a mock with six task templates and ten instances of each per test split, so the tasks are
not independent and the effects are tied to these templates; we cluster by task but do not model the
template, and a small effect such as P2's is absent in some templates. Only
Qwen3-8B carries confirmatory tests; Haiku was run with a single seed. The read-back is idealised: it
knows the true state before each write and the expected effect of every tool (Section~\ref{sec:method}).
The implausible variant of the wrong-read fault was added during development so that the contract had a
class of faults it could see; it is reported separately. Five further issues come from our own review of
the experiment. First, the grader counts every flag that is not on the faulted call as a false flag,
so the contract's correct flags on later re-reads of a corrupted record are counted as false; the
NOTICE wording is shared by all triggers and is not tuned to any of them. Second, the faulted call and
the fault's random draws come from the same seeded generator, so they are correlated within a (task,
fault, seed) triple; every arm still uses the same slot. Third, Qwen3-4B often fails to produce a
parseable reply, which depresses all its arms. Fourth, a fault only fires if the agent reaches an
eligible call, and arms change how often it does: on Qwen3-8B the wrong-read fault fired in a share
\RExpQwenEightAppliedNeverRead{} of \texttt{never} episodes but \RExpQwenEightAppliedAlwaysRead{} of
\texttt{always} episodes. Fifth, the step budget may cap repair: \RExpQwenEightStatusBudgetExhausted{}
Qwen3-8B and \RExpQwenFourStatusBudgetExhausted{} Qwen3-4B episodes end on it, and we did not test
whether a larger budget lets the small models repair wrong and partial writes. The uncertainty
threshold was chosen by rescoring development runs, where it would have fired far less often than it
did on the test split; we report it as run rather than tune it on the test split.
